\section{What the matched comparison resolves}
The direct dynamic--constant contrast removes the obvious static-four-logit versus MLP parameter-count difference from this new experiment, while retaining common feature pooling, initialization, draws, updates, and decoder paths. It strengthens the interpretation of a conditioning comparison under the specified bounded parameterization. It does not retrospectively show that earlier gains were capacity artifacts: the predecessor already included the same-sized A3 control, and this study changes the panel and residual formula.

The results distinguish three statements. First, the gate learns and responds to frames; the gradient and fixed-checkpoint variation records support that implementation statement. Second, a few annotations change from missed to matched, with small and seed-dependent net changes. Third, the practical criterion is not met: no primary seed obtains the minimum recall increase, and the positive pooled seed~1 effect accompanies negative equal-log recall and an additional failed frame. Evidence for the first two statements cannot be substituted for the third.

A scalar per layer aggregates detection gradients over a whole frame. The observed combination of recoveries and persistent failures is compatible with that architecture, but this experiment does not establish the mechanism causing any particular annotation to recover. Recoveries on already failed frames explain the arithmetic gap between count recall and AnyMiss; they do not show that the recovered objects are unimportant. A useful next hypothesis would require its own frozen design, new evidence, and explicit budget. No such expansion is executed here after failure of the current screen.
